% method-1.tex — Часть 1: Формирование гипотезы атаки
\documentclass[tikz, border=3mm]{standalone}
\usepackage{fontspec}
\usepackage[russian]{babel}
\setmainfont{PT Serif}
\setsansfont{PT Sans}
\usetikzlibrary{positioning, arrows.meta, shapes.geometric, shapes}

\begin{document}
\begin{tikzpicture}[
    node distance=1.0cm and 1.4cm,
    start/.style={ellipse, draw=black, fill=green!20, minimum width=1.6cm, minimum height=0.6cm, font=\scriptsize},
    process/.style={rectangle, draw=black, fill=blue!10, minimum width=2.3cm, minimum height=0.65cm, align=center, font=\scriptsize},
    decision/.style={diamond, draw=black, fill=orange!20, aspect=1.5, align=center, inner sep=1pt, font=\scriptsize},
    arrow/.style={->, >=Stealth, thick}
]

% Узлы
\node[start] (Start) {Начало};
\node[process, below=of Start] (P1) {Загрузка онтологии};
\node[process, below=of P1] (P2) {Построение эталонных метаграфов};
\node[process, below=of P2] (P3) {Подключение источников};
\node[process, below=of P3] (P4) {Нормализация событий};

\node[decision, below=of P4] (D1) {Событие\\известно?};

% ИЗМЕНЕНО: блок сдвинут чуть ниже и вправо
\node[process, below right=1.4cm and 2.6cm of D1] (P5) {Извлечение\\атрибутов};
\node[process, below=of P5] (P6) {Формирование\\$G_{\mathrm{obs}}$};

\node[decision, below=of D1] (D2) {Найден\\изоморфизм?};
\node[process, below=of D2] (P7) {Генерация\\гипотезы $H_i$};
\node[process, below=of P7] (P8) {Расчёт веса\\$W(H_i)$};
\node[decision, below=of P8] (D3) {$W > \theta$?};

% Архив
\node[process, right=2.4cm of D3, yshift=-0.8cm] (P9) {9. Архив\\ложного\\срабатывания};

% Стрелки
\draw[arrow] (Start) -- (P1);
\draw[arrow] (P1) -- (P2);
\draw[arrow] (P2) -- (P3);
\draw[arrow] (P3) -- (P4);
\draw[arrow] (P4) -- (D1);

% D1 → D2 и P5
\draw[arrow] (D1) -- node[left, font=\tiny] {да} (D2);
% ИЗМЕНЕНО: стрелка идёт вправо, затем поворачивает вниз и заходит строго сверху в P5.north
\draw[arrow] (D1.east) -- ++(1.5,0) |- node[above, pos=0.25, font=\tiny] {нет} (P5.north);
\draw[arrow] (P5) -- (P6);
\draw[arrow] (P6) -| (D2.east);

% Основной путь
\draw[arrow] (D2) -- node[left, font=\tiny] {да} (P7);
\draw[arrow] (P7) -- (P8);
\draw[arrow] (P8) -- (D3);

% Обход D2 → P9 (ложное срабатывание при отсутствии изоморфизма)
\coordinate[left=1.6cm of D2] (left_of_D2);
\coordinate[below=0.5cm of P9] (below_P9);
\draw[arrow] (D2.west) -- (left_of_D2) |- (below_P9) -- node[right, font=\tiny, pos=0.4] {нет} (P9.south);

% От D3 → P9 (вес ниже порога)
\draw[arrow] (D3.east) -- ++(0.6,0) |- node[right, font=\tiny, pos=0.25] {нет} (P9);

% Переход к части 2
\draw[arrow] (D3) -- ++(0,-0.8) node[below, font=\footnotesize] {(переход к части 2)};
\end{tikzpicture}
\end{document}